%This template is based on one provided by the American Physical Society for submission to its journals.

\documentclass[aps,twocolumn,showpacs,preprintnumbers]{revtex4}

%The following packages add LaTeX commands that make formatting and writing math easier

\usepackage{graphicx}  % Include figure files
\usepackage{subfigure}

\linespread{1.1}
\usepackage{fancyhdr}
\usepackage{parskip}
\usepackage[T1]{fontenc}
\usepackage{dcolumn}   % Align table columns on decimal point

\usepackage{bm}        % bold math
\usepackage{amsfonts}  % Common math fonts
\usepackage{amsmath}   % Common math functions
\usepackage{amssymb}   % Common math symbols

%The following custom commands simplify commonly used LaTeX commands

\newcommand{\pic}[2]{\begin{center} \includegraphics[scale=#1]{#2}\end{center}}
\newcommand{\re}[1]{\mathrm{Re}\left(#1\right)}
\newcommand{\im}[1]{\mathrm{Im}\left(#1\right)}
\newcommand{\bdot}[1]{\dot{ \bb {#1}}}
\newcommand{\bddot}[1]{\ddot{ \bb {#1}}}
\newcommand{\bidot}[1]{\dot{ \bi{ #1}}}
\newcommand{\biddot}[1]{\ddot{ \bi {#1}}}
\newcommand{\ep}{\varepsilon}
\newcommand{\for}{\quad \quad \mathrm{for} \quad\quad}
\newcommand{\then}{\quad \quad \implies \quad\quad}
\newcommand{\an}{\quad \quad \mathrm{and} \quad\quad}
\newcommand{\ifff}{\quad \quad \mathrm{if} \quad\quad}
\newcommand{\where}{\quad \quad \mathrm{where} \quad\quad}
\newcommand{\dg}{^\dagger}
\newcommand{\semi}{\quad \quad \mathrm{;} \quad\quad}
\newcommand{\paren}[1]{\left( #1 \right)}
\newcommand{\brac}[1]{\left[ #1 \right]}
\newcommand{\bra}[1]{\left\langle #1 \right|}
\newcommand{\exv}[1]{\left\langle #1 \right\rangle}
\newcommand{\pwisein}{\left\{ \begin{array}{ll}}
\newcommand{\pwiseout}{\end{array}\right.}
\newcommand{\ket}[1]{\left| #1 \right\rangle}
\newcommand{\bracket}[2]{\left\langle #1 | #2 \right\rangle}
\newcommand{\trace}[1]{\mathrm{Tr} \left( #1 \right)}
\renewcommand{\det}[1]{\mathrm{det}\left( #1 \right)}
\newcommand{\del}[1]{\frac{\partial}{\partial #1}}
\newcommand{\fulld}[1]{\frac{d}{d #1}}
\newcommand{\fulldd}[2]{\frac{d #1}{d #2}}
\newcommand{\dell}[2]{\frac{\partial #1}{\partial #2}}
\newcommand{\delltwo}[2]{\frac{\partial^2 #1}{\partial #2 ^2}}
\newcommand{\bb}{\mathbf}
\newcommand{\bi}{\boldsymbol}
\newcommand{\eq}[1]{\begin{equation} #1 \end{equation}}
\newcommand{\radhalf}{ \frac{ \sqrt{2}}{2}}
\newcommand{\sigx}{\left( \begin{array}{cc} 0 & 1\\ 1 & 0 \end{array}\right)}
\newcommand{\sigy}{\left( \begin{array}{cc} 0 & -i\\ i & 0 \end{array}\right)}
\newcommand{\sigz}{\left( \begin{array}{cc} 1 & 0\\ 0 & -1 \end{array}\right)}
\renewcommand{\matrix}[1]{\left( \begin{array} #1 \end{array}\right)}
\newcommand{\thermo}[3]{\left( \frac{\partial #1}{\partial #2} \right)_{#3}}
\newcommand{\coolfrac}[2]{\left( \frac{ #1}{ #2} \right)}

%\setlength{\parskip}{\baselineskip}
\setlength{\parindent}{10pt}

\begin{document}

\title{Your Title Here (e.g. Measurement of the speed of light using a laser distance meter)}

\author{Your Name$^\ast$}

\affiliation {\it Physics Department, University of California, Santa Barbara, CA 93106}

\date{submitted June 3, 2016}

\begin{abstract}  

%First sentence (or two) says something about the importance of the parameter being measured.
The speed of light is a fundamental physical constant with many counter-intuitive consequences. 
Its precise value is of use in many applications, including $\ldots$. 
%The next sentence (or two) summarizes the measurement approach used.
Using a laser distance meter and a rotating array of mirrors, we determined the speed of light from the change in reported distance as a function of the rotation frequency. 
%The final sentence (or two) state the result (measurement with its 68% confidence interval) and insight(s) concerning the measurement approach, if any.
We find light travels at $(3.1 \pm 0.2) \cdot 10^8$ m/s in air. 
The precision of our measurement was limited by uncertainty in spacing between mirrors in the array. 

\end{abstract}

%For fun, see PACS list at  https://www.aip.org/publishing/pacs/pacs-2010-regular-edition
\pacs{06.30.-k,42.25.Bs}

\maketitle 

%P1: Intro: scientific motivation

The notion that there exists an absolute limit on the speed at which light (or anything else, such as matter, or information) can move was introduced over a hundred years ago to explain...  \cite{EinsteinPHYSREV1905/7}. 
The fascinating, counter-intuitive consequences of this speed limit are collectively referred to as ``Special Relativity'' \cite{TaylorWheeler92}. 
In addition to satisfying curiosity, precise knowledge of the speed of light is important for a variety of technological applications, including... \cite{severalApllicationRefs}. 

%P2: More intro: approach, definitions, history, context, rationale

The first measurement of the speed of light, commonly referred to as $c$, was made in 1676 by Ole Roemer, who noted a seasonal variation in the periodicity of the phases of Jupiter's moon, Io \cite{BobisJAHH}.
The value obtained, $2.2 \cdot 10^8$ m/s, was limited in accuracy by errors in the accepted values for the diameters of Earth's and Jupiter's orbits \cite{refExample1}.
To date, the most precise measurement of the speed of light, commonly referred to as $c$, was made by researchers at the U.S. National Institute of Science and Technology (NIST) using... \cite{refExampleMANY}.

%P3: What you did, what you found 

Here we report a simple approach to measuring the speed of light in air using paper clips and bubble gum that yields a value of $(3.1 \pm 0.2) \cdot 10^8$ m/s.  
The precision of this measurement was limited by uncertainty in the spacing between mirrors.

%P4: Approach 1 - Setup (The nonsense below aims to give you a sense for style. It is purely imaginary.) 
Our approach was based on the one described by R. L. Goldberg \cite{GoldbergBOOK79}.
The light from a strobe lamp (ROX-ST1, Roxant) aimed at a fixed 3'' diameter front-surface mirror, was reflected onto a similar mirror, facing the first, attached to the end of a string to form a pendulum (Fig.~\ref{setup}).
The length of the pendulum thus created was measured with a laser distance meter and verified by tuning the frequency of the strobe lamp until the swinging mirror appeared to be stationary.
A detector (model 595391, Phantom YoYo), positioned behind the low point of the swinging mirror's arc, was triggered by the strobe with a variable delay.

%P5: Approach 2 - Methods (It may take several paragraphs to describe how you made your measurement.  Look at the paper from which you drew inspiration to get a sense of what to write.)
The delay corresponding to maximal intensity at the detector was recorded as a function of distance between the mirrors, over a range from 1 to 100 m.
The measurement was repeated five times for each configuration and 68\% confidence intervals were estimated from the standard deviation of these samples.

%P6: Approach 3 - Methods (There is no hard and fast rule about how many paragraphs to devote to describing your approach.  Use as many as you need in order to convey whatever information necessary for someone else to get the same result you got using your approach.)


%%%%%%%%%%%%%%%%%%%%Figure  1  begin %%%%%%%%%%%%%%%%%%%%%%%
\begin{figure}
%
\includegraphics[width=3.0in]{setup.jpg}
%
\caption{\label{setup} 
%Start with a clause summarizing what is presented in the figure.
Schematic of the mirror array and measurement setup. 
%Then add a sentence or two pointing out key features.
A wad of bubble gum anchors one mirror to the wall, while another mirror swings at the end of a pendulum three meters away.
A strobe illuminates the fixed mirror and a detector, positioned behind the mid-point of the swinging mirror's arc, is triggered a variable time after the strobe,}
%
\end{figure}
%%%%%%%%%%%%%%%%%%%%Figure 1 end %%%%%%%%%%%%%%%%%%%%%%%

 %%%%%%%%%%%%%%%%%%%%Figure  2  begin %%%%%%%%%%%%%%%%%%%%%%%
\begin{figure}[t]
%
\includegraphics[width=3.0in]{fakeplot.png}
%
\caption{\label{rawData} Fake data included for illustration of report layout. 
A properly formatted figure would have axis labels and error bars, of course, and a proper functional fit, instead of connected dots, if appropriate.
}
%
\end{figure}
%%%%%%%%%%%%%%%%%%%% Figure  2  end %%%%%%%%%%%%%%%%%%%%%%%

%P7: Results 1 - Start by describing the raw data.  Raw data is what you directly measured, with no calculations or other processing. 
An example of detected intensity as a function of delay time is plotted in Figure ~\ref{rawData}.  
The delay time corresponding to maximum intensity was determined for each mirror separation from a Gaussian fit to the data.
Uncertainty in the optimal delay time was derived from the least-squared-error (LSE) fit, taking uncertainties in both intensity and delay time into account.  

%P8: Results 2 - If you have two types of raw data, they may each require their own paragraph.  Oncer you are done describing the raw data, then describe the measurement itself, making it clear what, if any, calculations were done.
Optimal delay as a function of mirror separation is plotted in Figure~\ref{processedData}.
A LSE linear fit to the data yields a slope of $(4.0 \pm 0.9) \times 10^{-9}$ s/m.


%%%%%%%%%%%%%%%%%%%% Figure 3 begin %%%%%%%%%%%%%%%%%%%%%%%
\begin{figure}[!h]
%
\includegraphics[width=3.38in]{fakeLINEplot.png}
%
\caption{\label{processedData} More fake data included for illustration of report layout. 
A good caption starts with a descriptive clause, kind of like a title, followed by a sentence or two drawing the reader's attention to the ``take home message'' of the data displayed.
The interested reader will look to the main text to get complete details, so some of the information provided in this caption may be reiterated in the main text.}
%
\end{figure}
%%%%%%%%%%%%%%%%%%%% Figure 3 end %%%%%%%%%%%%%%%%%%%%%%%

%P9: Discussion - After presenting your result you may have ideas about it you want to share. You may want to say something about how your result compares with the results of others (so-called ``accepted values'', published elsewhere).

This result is consistent with measurements reported by G. Gedanken {\it et al.} \cite{GedankenAMJPHYS12}, who used a setup based on scotch tape and safety pins.  The approach used here afforded greater precision, $\pm 15\%$ compared to their $\pm 25\%$, although both managed to successfully avoid systematic errors imposed by reality. \cite{comment}

%P10:  Discussion - another good subject for discussion is your ideas for how the measurement approach could be improved.

%P11:  Discussion - use as many paragrahs as you like for the discussion.  It is fine to include matters of opinion, but be sure to support each assertion with clear reasoning. 


%P12:  Summary & Vision

In summary, we have measured the speed of light to be $(2.5 \pm 0.6) \times 10^8$ m/s.  
Our measurement suffered systematic errors due to the aging of the bubble gum used to anchor one mirror, and was limited in precision by the resolution of the photodetector.
The most marked improvements to this measurement approach would come from...

%P21: Acknowledgements

I thank my lab partner, Their Name, for insightful conversations and for proofreading this manuscript.  I am grateful to Alan Tran, Dillon Cislo and Ryan DeCrescent for patient explanations and thoughtful encouragement.  
This work was supported by Physics 25L lab fees.

\begin{thebibliography}{10}

\bibitem{EinsteinPHYSREV1905/7} A. Einstein, (From ``The Collected Papers, Vol 2, The Swiss Years: Writings, 1900–1909'', English Translation) ``On the Electrodynamics of Moving Bodies'', {\it Annalen Der Physik} {\bf 17}, 891–921 (1905); and ``On the Inertia of Energy Required by the Relativity Principle'', {\it Annalen Der Physik} {\bf 23}, 371–384 (1907).

\bibitem{TaylorWheeler92} E.F. Taylor \& J.A. Wheeler, {\it Spacetime Physics}, (WH Freeman, 1992).

\bibitem{severalApllicationRefs} I.~M. Author1 \&  I.~M. Author2, Article One's Title, {\it Abbrev. J. Title} {\bf vol}, pagei-pagef (year); A.~O. Scientist, Article Two's Title, {\it Abbrev. J. Title} {\bf vol}, pagei-pagef (year).

\bibitem{refExample1} O. Author, ``Article Title'', {\it Abbrev. J. Title} {\bf vol}, pagei-pagef (year).

\bibitem{refExample2} I.~M. Author1 \&  I.~M. Author2,  ``Article Title'', {\it Abbrev. J. Title} {\bf vol}, pagei-pagef (year).

\bibitem{refExampleMANY} T.~F. Author {\it et al.},  ``Article Title'', {\it Abbrev. J. Title} {\bf vol}, pagei-pagef (year).

\bibitem{BobisJAHH} L. Bobis \& J. Lequeux,  ``Cassini, R\o mer and the velocity of light'', {\it J. Astro. Hist. and Herit.}, {\bf 11}, 97-105 (2008).

\bibitem{GedankenAMJPHYS12} G. Gedanken {\it et al.}, ``Thought Experiments for Fun and Profit'', {\it Am. J. Phys.}, {\bf 0}, 3-14 (2012).

\bibitem{GoldbergBOOK79} R.~L. Goldberg, {\it The Best of Rube Goldberg}, (Prentice Hall, 1979). 

\bibitem{comment} It is ok to include footnotes among the references, like this.

\end{thebibliography}
\end{document}
